%%
%% This is file `example_pt.tex',
%% generated with the docstrip utility.
%%
%% The original source files were:
%%
%% coppe.dtx  (with options: `examplept')
%% 
%% This is one of the demonstration documents of the `coppe' document class:
%% the five main-language demos, the PDF/A build of the full example and the
%% side-by-side cover sheet.  It is generated from coppe.dtx by docstrip; edit
%% the .dtx source, not this file.
%% 
%% Copyright (C) 2008-2026 Vicente Helano F. Batista, George O. Ainsworth Jr.,
%% Geraldo Xexéo and Eduardo Mangeli. Released under the GNU General Public
%% License version 3 (see the COPYING.txt file).
%% 

\documentclass[dsc]{coppe}
\addbibresource{exemplo.bib}
\begin{document}
  \title{Demonstração da classe CoppeTeX 4.1 em português}
  \foreigntitle{CoppeTeX 4.1 class demonstration in English}
  \author{Nome do}{Autor}
  \advisor{Primeiro}{Orientador}{D.Sc.}{UFRJ}
  \examiner{Primeiro Examinador}{D.Sc.}{UFRJ}
  \examiner{Segundo Examinador}{D.Sc.}{UFRJ}
  \department{PESC}
  \date{05}{2026}
  \keyword{Multilíngue}
  \keyword{Demonstração}
  %% palavras-chave do resumo em idioma estrangeiro (3.1.2.1.5, Anexo F)
  \foreignkeyword{Multilingual}
  \foreignkeyword{Demonstration}
  \maketitle
  \frontmatter
  \begin{abstract}
    Este é o resumo principal em português, idioma escolhido por padrão
    quando nenhuma opção de idioma é passada para a classe \texttt{coppe}.
    Lorem ipsum dolor sit amet, consectetur adipiscing elit.
  \end{abstract}
  \begin{foreignabstract}
    This is the foreign abstract in English. With a Brazilian-Portuguese
    main language, \texttt{foreignabstract} defaults to English by class
    convention.
  \end{foreignabstract}
  % Listas pré-textuais: só o sumário é obrigatório (Manual 3.1.2.1.6).
  % As listas de figuras, tabelas, abreviaturas e siglas e símbolos são
  % opcionais (3.1.2.2); as de quadros, programas e algoritmos entram quando
  % o trabalho tem um deles (Norma COPPE, §10). Esta demonstração não tem
  % ilustrações, por isso só imprime o sumário.
  \tableofcontents
  \mainmatter
  \chapter{Introdução}
  Este capítulo demonstra que o corpo do trabalho é composto em português,
  com legendas, citações e referências automaticamente em português.
  \section{O arquivo de depósito: PDF/A}
  O depósito na UFRJ é exclusivamente digital desde a Resolução CEPG
  n.~246/2023, e o Manual (2.2d) exige o arquivo final em \textbf{PDF/A}, o
  formato de arquivamento de longo prazo. PDF/A é o PDF restrito pela norma
  ISO~19005 para preservação: todas as fontes vão embutidas, as cores declaram
  um perfil (sRGB) que as torna independentes do equipamento, os metadados vão
  num bloco XMP que os repositórios leem, e nada depende de fora --- sem
  conteúdo externo, JavaScript ou criptografia. Assim o arquivo abre igual daqui
  a décadas, em qualquer computador.

  A classe usa a variante \textbf{PDF/A-2b}: a parte 2 admite transparência,
  que a parte 1 proíbe e que caixas coloridas e gráficos produzem, e o nível
  \emph{b} garante a reprodução visual fiel. Basta a opção \texttt{pdfa} em
  \verb|\documentclass[dsc,pdfa]{coppe}|: a classe carrega o \texttt{pdfx},
  monta os metadados a partir do título, do autor e das palavras-chave, e
  embute o perfil de cor. Ligue-a cedo, como faz o \texttt{max-exemplo.tex}: um PDF
  ou uma imagem incluídos fora da norma aparecem no dia em que entram, e não na
  véspera do depósito. Para conferir, use o validador veraPDF, perfil PDF/A-2b.

  \section{Fluxos de escrita}
  O \hologo{pdfLaTeX} mantém no máximo dezesseis arquivos abertos para escrita,
  e cada lista, índice ou glossário ocupa um. Por isso a classe carrega o
  \texttt{morewrites}. A opção \texttt{semmorewrites} o desliga e poupa cerca de
  meio segundo por passada, mas só serve a trabalho sem índice remissivo, sem
  vários índices e sem glossário.
\end{document}
%% 
%%
%% End of file `example_pt.tex'.
